\documentclass{article}

\usepackage[T1]{fontenc}

\usepackage{booktabs,graphicx,siunitx}

\begin{document}

\section{Exploratory factor analysis (EFA)}

\begin{tabular}{llll}
\toprule
KMO & Bartlett's $\chi$\textsuperscript{2} & df & p \\
\midrule
.83 & 2117.6 & 55 & $<$.001 \\
\bottomrule
\end{tabular}

\begin{tabular}{llll}
\toprule
Factor & Eigenvalue & \% variance & Cumulative \% \\
\midrule
F1 & 2.62 & 23.8 & 23.8 \\
F2 & 2.40 & 21.8 & 45.6 \\
F3 & 1.94 & 17.7 & 63.3 \\
\bottomrule
\end{tabular}

\begin{tabular}{lllll}
\toprule
Item & F1 & F2 & F3 & Communality \\
\midrule
esg1 &  & .83 &  & .68 \\
esg2 &  & .76 &  & .57 \\
esg3 &  & .79 &  & .62 \\
esg4 &  & .71 &  & .53 \\
norm1 & .84 &  &  & .71 \\
norm2 & .77 &  &  & .60 \\
norm3 & .81 &  &  & .65 \\
norm4 & .80 &  &  & .65 \\
intent1 &  &  & .84 & .70 \\
intent2 &  &  & .79 & .59 \\
intent3 &  &  & .78 & .65 \\
\bottomrule
\end{tabular}

\begin{tabular}{lccc}
\toprule
 & F1 & F2 & F3 \\
\midrule
F1 & 1.00 &  &  \\
F2 & $-$0.11 & 1.00 &  \\
F3 & 0.25 & 0.37 & 1.00 \\
\bottomrule
\end{tabular}

\includegraphics{figures/01_efa/figure_scree.png}

Check KMO ($>$ .60 acceptable, $>$ .70 preferred; Kaiser \& Rice, 1974) and a significant Bartlett's test (p $<$ .001) for suitability. Parallel analysis (Horn, 1965) is the recommended retention rule - it corrects the upward bias of the Kaiser eigenvalue $>$ 1 rule (Zwick \& Velicer, 1986). Table 3 shows the rotated pattern matrix: each loading is the unique contribution of the factor to that item. Loadings |$<$ .32| (\textasciitilde{}10\% variance) are suppressed; a loading $\ge$ .32 on two or more factors flags a cross-loader (Tabachnick \& Fidell). For oblique rotation (oblimin), the Φ matrix shows interfactor correlations - if all |Φ| $<$ .32, orthogonal rotation (varimax) may be preferred. EFA is exploratory - it discovers structure; to confirm a proposed factor structure run a CFA, ideally on a separate sample (Watkins, 2018).

EFA (KMO = .83, Bartlett's $\chi$\textsuperscript{2}(55) = 2117.6, p $<$ .001) with parallel analysis retained 3 factors explaining 63.3\% of variance (oblimin rotation).

Statistical basis: KMO sampling-adequacy verbal labels ($\ge$ .60 acceptable, $\ge$ .70 preferred) (Kaiser, H. F., Rice, J. (1974). "Little Jiffy, Mark IV." Educational and Psychological Measurement, 34(1), 111-117.); Parallel analysis as the preferred factor-retention rule (Horn, J. L. (1965). "A rationale and test for the number of factors in factor analysis." Psychometrika, 30(2), 179-185.); Kaiser eigenvalue $>$ 1 rule systematically over-extracts (Zwick, W. R., Velicer, W. F. (1986). "Comparison of five rules for determining the number of components to retain." Psychological Bulletin, 99(3), 432-442.); EFA best-practice reporting standard (Watkins, M. W. (2018). "Exploratory Factor Analysis: A Guide to Best Practice." Journal of Black Psychology, 44(3), 219-246.).

\end{document}